%% kapa: exam=1 round=36 year=2025 period=2 form=A subject=appraisal_law subject_ko=감정평가관계법규 seq=10 number=10
%% source: Q-Net 감정평가사 기출문제 (cst00309) artlSeq=5231525 file=raw/1008/5231525/2025년 제36회 감정평가사 1차시험 2교시 문제지 원본.pdf
%% conversion: kordoc 4.24.51 -> tools/gen1.py (auto); review=proofread-sonnet; flags=none
\begin{question}{10}
국토의 계획 및 이용에 관한 법령상 기반시설부담구역에 관한 설명으로 옳지 않은 것은?
\begin{choices}
  \item 시장은 기반시설부담구역을 지정 또는 변경하려면 주민의 의견을 들어야 한다.
  \item 기반시설부담구역의 지정고시일부터 1년이 되는 날까지 기반시설설치계획을 수립하지 아니하면 그 1년이 되는 날의 다음날에 기반시설부담구역의 지정은 해제된 것으로 본다.
  \item 시장은 기반시설설치비용의 관리 및 운용을 위하여 기반시설부담구역별로 특별회계를 설치하여야 한다.
  \item 기반시설부담구역에서 150제곱미터의 단독주택으로 증축하는 행위는 기반시설설치비용의 부과대상이다.
  \item 공원은 기반시설부담구역에 설치가 필요한 기반시설에 해당한다.
\end{choices}
\end{question}
\answer{4}
